\section{Safety-Analyse (Kap. 06)}
\label{sec:safety}

\subsection{Angriffsvektoren und Guardrails}

\textbf{Indirekte Prompt-Injection aus Lernmaterialien.} Ein Angreifer kann in
einem hochgeladenen PDF Anweisungen platzieren, die beim semantischen Abruf in
den Tutor-Prompt gelangen. Genau diese Vermischung von Daten und Instruktionen
ermöglicht indirekte Prompt-Injection bis hin zur Manipulation nachgelagerter
API-Aufrufe \parencite{greshakeEtAl2023}. StudyMummy filtert einige direkte
Muster in Nutzernachrichten, isoliert RAG nach Nutzer und Dokument und markiert
Dokumenttext im Systemprompt als nicht vertrauenswürdig. Der Reviewer soll
befolgte Dokumentanweisungen erkennen. Diese Guardrails reduzieren das Risiko,
sind aber nicht vollständig: Chunks und \texttt{extra\_context} werden nicht
semantisch geprüft, und Review findet erst nach Toolausführung statt.

Passend wären mehrere Grenzen: Dokumente und Frontend-Kontext immer als
unvertrauenswürdige, zitierbare Datensätze serialisieren, Retrieval-Ergebnisse
mit Provenienz versehen, geplante Toolargumente vor Ausführung gegen eine
deterministische Policy prüfen, Dokumentkontext darf niemals neue Tools oder
Identitäten bestimmen. Ergänzend sollten Injection-Testkorpora, Rate Limits
und Quarantäne bei verdächtigen Chunks eingesetzt werden. Ein weiterer
LLM-Filter allein wäre wegen korrelierter Fehlentscheidungen kein verlässlicher
Sicherheitsanker.

\textbf{Toolmissbrauch und ungewollte Seiteneffekte.} Direkte Aufforderungen
könnten Coins erzeugen oder ohne ausreichende Bestätigung Termine anlegen.
Vorhanden sind aktionsabhängige Toolrechte, eine zweite Scope-Prüfung im
LLM-Adapter, gebundene Nutzeridentität sowie beim Belohnungstool Betragsgrenze,
Eigentumsprüfung, Zeilensperre und Doppelbelohnungsschutz. Die einfache Suche
nach Wörtern wie \enquote{coin} ist umgehbar und kann legitime Nachrichten
blockieren. Beim Kalender fehlen Nutzerbestätigung, Idempotenzschlüssel und die
fachliche Prüfung, dass Ende nach Beginn liegt. Da der Reviewer erst nach der
Ausführung arbeitet, kann er diese Wirkung nicht verhindern.

Ein Pre-Execution-Monitor sollte daher Tool, Argumente, Identität,
Wertebereiche und erwartete Aktion prüfen. Schreibende Aufrufe werden zunächst
als \texttt{pending} protokolliert, Kalendertermine werden erst nach
Bestätigung committet. Für Belohnungen reichen wegen geringer Reversibilität
enge automatisierte Grenzen und Auditlog. Verbleibende Risiken sind
LLM-Nichtdeterminismus, Promptvarianten außerhalb der Filter, vergiftetes
Langzeitgedächtnis, Kostenmissbrauch durch große Uploads sowie personenbezogene
Inhalte in Logs und Vektorspeicher.

\subsection{Human-in-the-Loop}

Vollständige Freigabe jedes Tutor-Turns würde die Lerninteraktion unbrauchbar
verlangsamen. Geeignet ist ein risikobasiertes Human-in-the-Loop: Lesen,
Retrieval, Antwortbewertung und Textentwürfe laufen autonom, ein sichtbarer
Bestätigungsdialog zeigt bei persistenten Seiteneffekten Toolname, Argumente,
Quelle und erwartete Wirkung. Die lernende Person bestätigt Kalendertermine
und kann gespeicherte Erinnerungen löschen. Lehrende oder Betreiber prüfen
stichprobenartig abgelehnte Antworten und Injection-Traces. HITL begrenzt den
Schadensradius, ersetzt aber weder Rechteprüfung noch Monitoring: Menschen
können unter Zeitdruck routinemäßig bestätigen und sehen versteckte
Kontextmanipulation nicht. Deshalb muss die technische Policy auch bei
menschlicher Freigabe unverändert gelten.
